% !TEX root = ../линейная_алгебра.tex

\section{Ранг матрицы над полем}

\subsection{Ранг набора векторов и ранг матрицы}

\begin{Definition}{Ранг набора векторов}{}
    \emph{Рангом} набора векторов $v_1, \ldots, v_n$ называется размерность их линейной оболочки:
    \[
        \mathrm{rk}(v_1, \ldots, v_n) = \dim \mathrm{Lin}(v_1, \ldots, v_n).
    \]
\end{Definition}

Пусть $k$ — поле, $A \in k^{m \times n}$.

\begin{Definition}{Столбцовый и строковый ранги}{}
    \emph{Столбцовым рангом} матрицы $A$ называется ранг системы её столбцов:
    \[
        \mathrm{rk}_{\mathrm{col}}(A) = \mathrm{rk}(A[{:},1], \ldots, A[{:},n]).
    \]
    \emph{Строковым рангом} матрицы $A$ называется ранг системы её строк:
    \[
        \mathrm{rk}_{\mathrm{row}}(A) = \mathrm{rk}(A[1,{:}], \ldots, A[m,{:}]).
    \]
\end{Definition}

\subsection{Инвариантность рангов при элементарных преобразованиях строк}

\begin{Statement}{Инвариантность рангов}{}
    Столбцовый и строковый ранги матрицы $A$ не меняются при элементарных преобразованиях строк матрицы $A$.
\end{Statement}

\begin{proof}
Любое элементарное преобразование строк задаётся умножением слева на обратимую матрицу $U \in M_m(k)$: то есть $A' = UA$.

\textbf{Инвариантность столбцового ранга.} Покажем, что $\mathrm{rk}_{\mathrm{col}}(A') = \mathrm{rk}_{\mathrm{col}}(A)$, то есть что из линейной зависимости столбцов $A'$ следует линейная зависимость столбцов $A$ и наоборот.

Пусть $\lambda_1 A'[{:},i_1] + \cdots + \lambda_\ell A'[{:},i_\ell] = 0$. Так как $A'[{:},j] = UA[{:},j]$, получаем $U(\lambda_1 A[{:},i_1] + \cdots + \lambda_\ell A[{:},i_\ell]) = 0$. Поскольку $U$ обратима, $\lambda_1 A[{:},i_1] + \cdots + \lambda_\ell A[{:},i_\ell] = 0$. Обратное рассуждение симметрично (с $U^{-1}$). Следовательно, $\mathrm{rk}_{\mathrm{col}}(A') = \mathrm{rk}_{\mathrm{col}}(A)$.

\textbf{Инвариантность строкового ранга.} Элементарные преобразования строк — это именно перестановки строк, умножение строки на скаляр и прибавление одной строки к другой. Каждое из них не меняет линейную оболочку строк, откуда $\mathrm{rk}_{\mathrm{row}}(A') = \mathrm{rk}_{\mathrm{row}}(A)$.
\end{proof}

\begin{Remark}{}{}
    Если матрица $A$ элементарными преобразованиями строк приведена к ступенчатому виду с $r$ ненулевыми строками (окаймлённая единичная матрица $\begin{pmatrix} E_r & 0 \\ 0 & 0 \end{pmatrix}$), то количество единиц на диагонали равно $r = \mathrm{rk}(A)$.
\end{Remark}

\subsection{Равенство столбцового и строкового рангов}

\begin{Theorem}{Равенство рангов}{}
    Для любой матрицы $A \in k^{m \times n}$:
    \[
        \mathrm{rk}_{\mathrm{col}}(A) = \mathrm{rk}_{\mathrm{row}}(A).
    \]
    Общую величину называют просто \emph{рангом} матрицы и обозначают $\mathrm{rk}(A)$.
\end{Theorem}

\begin{proof}
Элементарными преобразованиями строк приведём $A$ к виду $D = \begin{pmatrix} E_r & 0 \\ 0 & 0 \end{pmatrix}$, где $r$ — число ненулевых строк. По доказанному выше, при этом оба ранга не изменились. Для матрицы $D$ очевидно: $\mathrm{rk}_{\mathrm{col}}(D) = r$ (ровно $r$ ненулевых столбцов, линейно независимых) и $\mathrm{rk}_{\mathrm{row}}(D) = r$ (ровно $r$ ненулевых строк, линейно независимых).
\end{proof}

\subsection{Минорный ранг}

\begin{Definition}{Минор}{}
    \emph{Минором} порядка $s$ матрицы $A \in k^{m \times n}$ называется определитель подматрицы $A[\{i_1,\ldots,i_s\},\{j_1,\ldots,j_s\}]$, образованной $s$ строками и $s$ столбцами.
\end{Definition}

\begin{Definition}{Минорный ранг}{}
    \emph{Минорным рангом} матрицы $A$ называется максимальный порядок ненулевого минора:
    \[
        \mathrm{rk}_{\mathrm{min}}(A) = \max\{s \mid \exists \text{ ненулевой минор порядка } s\}.
    \]
\end{Definition}

\begin{Theorem}{Все ранги совпадают}{}
    Для любой матрицы $A \in k^{m \times n}$:
    \[
        \mathrm{rk}_{\mathrm{col}}(A) = \mathrm{rk}_{\mathrm{row}}(A) = \mathrm{rk}_{\mathrm{min}}(A).
    \]
\end{Theorem}

\begin{proof}
Обозначим $r = \mathrm{rk}(A)$. Покажем два неравенства.

\textbf{$\mathrm{rk}_{\mathrm{min}}(A) \leq r$.} Рассмотрим произвольный минор порядка $s > r$. Он образован $s$ столбцами, которые линейно зависимы (так как $\mathrm{rk}_{\mathrm{col}}(A) = r < s$), значит определитель равен нулю.

\textbf{$\mathrm{rk}_{\mathrm{min}}(A) \geq r$.} Существуют $r$ линейно независимых столбцов с индексами $j_1, \ldots, j_r$. Рассмотрим подматрицу $B = A[{:},\{j_1,\ldots,j_r\}]$ — она имеет ранг $r$, то есть содержит $r$ линейно независимых строк с индексами $i_1, \ldots, i_r$. Тогда подматрица $A[\{i_1,\ldots,i_r\},\{j_1,\ldots,j_r\}]$ — квадратная, её строки линейно независимы, значит её определитель ненулевой. Это минор порядка $r$.
\end{proof}

\subsection{Тензорный ранг}

\begin{Definition}{Тензорный ранг}{}
    \emph{Тензорный ранг} матрицы $A \in k^{m \times n}$ определяется следующим образом:
    \begin{itemize}
        \item $\mathrm{rk}_{\otimes}(A) = 0$, если $A = 0$;
        \item $\mathrm{rk}_{\otimes}(A) = 1$, если $A \neq 0$ и существуют векторы
              $\alpha \in k^m$, $\beta \in k^n$ такие, что $A_{ij} = \alpha_i \beta_j$;
        \item в общем случае $\mathrm{rk}_{\otimes}(A) = \min\bigl\{r \mid A = A_1 + \cdots + A_r,\;
              \mathrm{rk}_{\otimes}(A_i) = 1\bigr\}$.
    \end{itemize}
\end{Definition}
\begin{Theorem}{Тензорный ранг равен рангу}{}
    $\mathrm{rk}_{\otimes}(A) = \mathrm{rk}(A)$.
\end{Theorem}

\begin{proof}
\textbf{$\mathrm{rk}_{\otimes}(A) \leq \mathrm{rk}(A)$.} Пусть $\mathrm{rk}(A) = r$. Приведём $A$ к виду $D = \begin{pmatrix} E_r & 0 \\ 0 & 0 \end{pmatrix}$ (умножением на обратимые). Матрица $D$ является суммой $r$ матриц ранга $1$ (по одной на каждую единицу диагонали), откуда $\mathrm{rk}_{\otimes}(A) = \mathrm{rk}_{\otimes}(D) \leq r$.

\textbf{$\mathrm{rk}_{\otimes}(A) \geq \mathrm{rk}(A)$.} Пусть теперь $r = \mathrm{rk}_{\otimes}(A)$, то есть $A = A_1 + \cdots + A_r$, где $\mathrm{rk}(A_i) = 1$ и $r$ минимально. Каждый столбец суммы есть сумма соответствующих столбцов слагаемых, поэтому $\Lin(A) \subseteq \Lin(A_1) + \cdots + \Lin(A_r)$, и так как размерность суммы подпространств не превосходит суммы размерностей:
\[
    \mathrm{rk}(A) = \mathrm{rk}(A_1 + \cdots + A_r) \leq \mathrm{rk}(A_1) + \cdots + \mathrm{rk}(A_r) = r.
\]
Следовательно, $r \geq \mathrm{rk}(A)$, то есть $\mathrm{rk}_{\otimes}(A) \geq \mathrm{rk}(A)$.
\end{proof}

\subsection{Неравенства для рангов}

\begin{Theorem}{Неравенства для рангов}{}
    Для матриц $A \in k^{m \times n}$, $B \in k^{n \times p}$:
    \begin{enumerate}
        \item $\mathrm{rk}(A + B) \leq \mathrm{rk}(A) + \mathrm{rk}(B)$.
        \item $\mathrm{rk}(AB) \leq \min(\mathrm{rk}(A),\, \mathrm{rk}(B))$.
        \item Если $U \in GL_m(k)$, то $\mathrm{rk}(UA) = \mathrm{rk}(A)$. Если $U \in GL_n(k)$, то $\mathrm{rk}(AU) = \mathrm{rk}(A)$.
    \end{enumerate}
\end{Theorem}

\begin{proof}
Пункт 1 следует из $\mathrm{Lin}(A+B) \subseteq \mathrm{Lin}(A) + \mathrm{Lin}(B)$.

Для пункта 2: столбцы $AB$ лежат в линейной оболочке столбцов $A$, откуда $\mathrm{rk}(AB) \leq \mathrm{rk}(A)$. Оценка $\mathrm{rk}(AB) \leq \mathrm{rk}(B)$ следует транспонированием: $\mathrm{rk}(AB) = \mathrm{rk}(B^T A^T) \leq \mathrm{rk}(B^T) = \mathrm{rk}(B)$.

Пункт 3 — следствие пункта 2: $\mathrm{rk}(UA) \leq \mathrm{rk}(A)$ и $\mathrm{rk}(A) = \mathrm{rk}(U^{-1} \cdot UA) \leq \mathrm{rk}(UA)$.
\end{proof}

\subsection{Теорема Кронекера--Капелли}

\begin{Theorem}{Кронекера--Капелли}{}
    Система линейных уравнений $Ax = b$, где $A \in k^{m \times n}$, совместна тогда и только тогда, когда
    \[
        \mathrm{rk}(A \mid b) = \mathrm{rk}(A).
    \]
\end{Theorem}

\begin{proof}
Система совместна тогда и только тогда, когда $b$ выражается через столбцы $A$, то есть $b \in \mathrm{Lin}(A[{:},1], \ldots, A[{:},n])$. Это равносильно тому, что добавление $b$ к системе столбцов не увеличивает размерность линейной оболочки, то есть $\mathrm{rk}(A \mid b) = \mathrm{rk}(A)$.
\end{proof}

\subsection{Прямые суммы подпространств}

\begin{Definition}{Внутренняя прямая сумма}{}
    Говорят, что $V$ \emph{раскладывается во внутреннюю прямую сумму} подпространств $W_1, \ldots, W_k \leq V$, и пишут $V = W_1 \oplus \cdots \oplus W_k$, если для каждого $v \in V$ существует единственное представление $v = w_1 + \cdots + w_k$, $w_i \in W_i$.
\end{Definition}

\begin{Remark}{}{}
    Пусть $e_1, \ldots, e_k$ — ненулевые элементы $V$. Тогда $e_1, \ldots, e_k$ — базис $V$ тогда и только тогда, когда
    \[
        V = \mathrm{Lin}(e_1) \oplus \cdots \oplus \mathrm{Lin}(e_k).
    \]
\end{Remark}

\begin{Theorem}{Критерий для двух слагаемых}{}
    Для подпространств $W_1, W_2 \leq V$:
    \[
        V = W_1 \oplus W_2 \iff V = W_1 + W_2 \text{ и } W_1 \cap W_2 = \{0\}.
    \]
\end{Theorem}

\begin{proof}
$(\Rightarrow)$: Если $w \in W_1 \cap W_2$ и $w \neq 0$, то $w = w + 0 = 0 + w$ — два различных разложения нуля, противоречие.

$(\Leftarrow)$: Если $v = w_1 + w_2 = w_1' + w_2'$, то $w_1 - w_1' = w_2' - w_2 \in W_1 \cap W_2 = \{0\}$, значит $w_1 = w_1'$, $w_2 = w_2'$.
\end{proof}

\begin{Theorem}{Критерий для нескольких слагаемых}{}
    $V = W_1 \oplus \cdots \oplus W_k$ тогда и только тогда, когда:
    \begin{enumerate}
        \item $W_1 + \cdots + W_k = V$;
        \item для каждого $j$: $W_j \cap (W_1 + \cdots + \widehat{W_j} + \cdots + W_k) = \{0\}$.
    \end{enumerate}
\end{Theorem}

\begin{Theorem}{Базис прямой суммы}{}
    Пусть $W_1, \ldots, W_k \leq V$, и пусть $e_{i,1}, \ldots, e_{i,d_i}$ — базис $W_i$. Тогда
    \[
        V = W_1 \oplus \cdots \oplus W_k \iff \{e_{i,j} \mid 1 \leq i \leq k,\; 1 \leq j \leq d_i\} \text{ — базис } V.
    \]
\end{Theorem}

\begin{proof}
$(\Rightarrow)$: Любой $v \in V$ единственным образом раскладывается как $v = \sum_i w_i$, $w_i \in W_i$, а каждое $w_i$ единственным образом раскладывается по базису $W_i$. Таким образом, $\{e_{i,j}\}$ — базис $V$.

$(\Leftarrow)$: Из того, что $\{e_{i,j}\}$ — базис $V$, следует, что любой вектор единственным образом раскладывается по ним, что равносильно единственности разложения $v = \sum_i w_i$ с $w_i \in W_i$.
\end{proof}